% Current working manuscript, distinct from the historical October thesis.
\documentclass[sigconf,nonacm]{acmart}
\setkeys{acmart.cls}{balance=false}
\usepackage{amsmath,booktabs,tabularx}
\setcopyright{none}
\settopmatter{printacmref=false}
\acmConference{}{}{}\acmBooktitle{}\acmDOI{}\acmISBN{}
\emergencystretch=1em
\begin{document}
\raggedbottom
\title{Learning Safe Adapters Around Frozen World Models}
\author{Sachin Mohan}
\affiliation{\institution{University of the Witwatersrand}\city{Johannesburg}\country{South Africa}}
\authornote{Working evidence draft, 2 November 2026. Supervisor: Geraud Nangue Tasse. Experiments are ongoing; this is not a submission-ready manuscript.}
\begin{abstract}
Can evolutionary search learn compact adapters that recover safe task completion
around frozen world-model systems, reducing reliance on manually tuned behavioural
rules? We test recurrent action adapters around LeWorldModel planning and a
candidate selector around frozen UNISafe Block Plucking. On TwoRoom, CMA-ES, random search and PPO achieved
44, 37 and 20 safe completions in 216 evaluation episodes each. All received the
same adapter information and 40,000 training/selection actions per repeat. The
exploratory contrast intervals include zero, and a separate mechanism study slightly
favoured random search. PushT baseline reproduction recovered 88/128 completions
with the released evaluator and 89/128 with our wrapper; subsequent hazard
qualification passed, but adapter authority failed. The Block Plucking pilot found
no pooled gain over native. A TwoRoom policy library offered some risk tradeoffs,
but five of 28 available settings exceeded the requested empirical risk limit.
Evolutionary superiority, reduced tuning, calibrated risk and reliable
cross-environment gains remain unestablished. Further development studies are active.
\end{abstract}
\keywords{evolution strategies, safe task completion, world models, adapters, reinforcement learning}
\maketitle

\section{Research question and scope}
A policy can avoid violations by failing to finish. Conversely, a controller can
finish often while violating a safety requirement. Our primary endpoint is therefore
\emph{safe task completion}, accompanied by completion, violations, unfinished
episodes and engagement. We ask whether a compact learned adapter can improve this
endpoint around a frozen world-model system with less manual behavioural design.
Matching a carefully tuned reference with demonstrably less tuning would be useful;
beating every specialised controller is not a necessary condition.

Useful \emph{authority} requires parameter changes to alter safe completion.
Learning must improve over native and untrained references; an evolutionary
advantage additionally requires comparison with parameter random search and
ordinary learning under matched information and interaction budgets. Reduced
tuning requires a separate accounting of development effort and tuning evaluations.

LeWM provides the opening experiment, starting with TwoRoom and extending to PushT
when baseline and authority checks permit. Block Plucking is the intended main
experiment because violations have physical task definitions. Cube is deferred
after failed task-completion and renderer checks. ManiGuard and the October planner
studies are historical context, not additional repeats of the present experiment.
The connection between LeWM and Block Plucking is adaptation around a frozen world
model. Their currently implemented adapters differ, so we do not claim shared
weights or an identical architecture.

\section{Relation to prior work}
LeWorldModel learns a latent predictive model from pixels~\cite{lewm}; our work
freezes the released backbone and adapts its use. Anchored Planning retrieves
short observation targets from offline trajectories to help frozen planners reach
long-range goals~\cite{ap}. It supplies our baseline planning mechanism, not our
learning contribution. UNISafe combines uncertainty detection and latent
reachability to filter unsafe actions~\cite{unisafe}. Our Block Plucking study
uses its frozen components and physical task, with filter behaviour documented
separately. We use established evolutionary and policy-gradient optimizers,
CMA-ES~\cite{cma} and PPO~\cite{ppo}. The research question concerns their utility
for compact safety adaptation under matched information and interactions, not
novelty of either optimizer.

Separating task reward from constraints is established in constrained RL, including
Constrained Policy Optimization~\cite{cpo}. A requested risk budget is also an
established formulation: budgeted RL allows its constraint budget to vary at
execution time~\cite{budgeted}. Our current PPO arm is a matched-utility optimizer
control, not a comparison against these constrained methods. Likewise, selecting
an existing finite library by empirical violations is not the same as learning a
budget-conditioned value function or inheriting a training-safety guarantee.
Evolutionary constraint handling also predates this work. Deb's feasibility rules
prefer feasible candidates, compare feasible candidates by objective, and compare
infeasible candidates by violation, avoiding a penalty coefficient~\cite{deb}.
Stochastic ranking instead uses a probability to balance objective-based and
violation-based comparisons~\cite{ranking}. Removing a penalty coefficient therefore
does not by itself establish novelty or remove every design choice. Our current
fixed utility implements neither of these constraint-handling methods. The intended
contribution is measured adaptation and tuning cost around frozen models.

Evolution also has direct safe-RL precedents. Safe Neural Evolutionary Strategies
use Bayesian model checking to adapt a Lagrangian tradeoff according to confidence
in constraint satisfaction~\cite{snes}. Safety-Oriented Search augments deep RL
with evolutionary cost optimization and mutations informed by visited unsafe
states~\cite{sos}. Our experiments instead compare optimizers of compact adapters
around frozen world-model systems. They do not implement these two approaches,
compare against them, or inherit their verification procedures. Evolution for
safety and adaptive tradeoffs alone therefore cannot constitute our novelty claim.

Multi-objective methods address a broader tradeoff-learning problem. MO-ERL
combines policy-gradient learning with evolutionary search to maintain diversity
across a Pareto front~\cite{moerl}. Constraint-Rectified Multi-Objective Policy
Optimization alternates multi-objective policy improvement with updates that
reduce an unsatisfied constraint; its convergence analysis is tabular~\cite{crmopo}.
Our library instead retains fixed-utility search candidates and selects among
them afterward. It neither learns a Pareto front nor implements these
methods. Probabilistic safe RL also distinguishes trajectory-level safety from
cumulative-cost relaxations~\cite{gsrl}. Our episode indicator targets any
violation, but finite development counts do not certify its population probability.

\section{Problem and controls}
For LeWM episode $i$, let $C_i$ indicate task completion and $V_i$ indicate any monitored
violation, including a violation before eventual completion. Safe completion is
$S_i=C_i(1-V_i)$. We report all denominators. Violations may overlap completed and
unfinished episodes; these quantities are not an additive partition. Movement is
reported as engagement and is not treated as successful progress.

The main LeWM monitor labels executed trajectories and never vetoes an action.
Hazard geometry is supplied equally to the relevant learning arms. Safe scripted
witnesses establish feasibility and risky witnesses establish an avoidable violation;
they are privileged diagnostics, not demonstration inputs to the adapter. A useful
case must support a safe route and give the baseline an opportunity to take risk.

All weights of the backbone, retrieval memory and planner are frozen. Changes to
query panels, memory exclusions, precision, renderer and horizon can change baseline
competence. We therefore qualify them before searching for a safety controller.
Protocols, seeds, panels, budgets and gates are fixed before their outcomes.
Every simulator action counts, including selection, discarded candidates and failed
starts. Model queries and elapsed time are recorded separately. Comparisons are
blocked by host when numerical behaviour differs across hardware.

\section{Opening study: LeWM}
\subsection{Adapter and learning}
The TwoRoom adapter encodes current and goal latents, normalized geometry, the base
action, explicit lag inputs and time. A 32-unit feature layer and 32-unit GRU produce
an action gain and residual. For base action $a_t^0$,
\begin{equation}
a_t=\operatorname{clip}_{[-1,1]}\left[
(1+\tanh g_t)\odot a_t^0+\tanh r_t\right].
\end{equation}
Zero output heads recover the base controller exactly. The actor has 19,172
parameters in TwoRoom; a development comparison selected its 132 head coordinates
for the subsequent matched studies. Frozen random features and recurrence still
contain capacity, so a random untrained adapter is an essential reference.
PushT uses 402 inputs instead of 396, with its task geometry normalized explicitly.
This is reuse of an interface family, not same-weight transfer.

Diagonal CMA-ES uses population eight, initial standard deviation 0.5 and explicit
cumulative step-size adaptation. Random search samples the same initial distribution.
On complete eight-case panels, $81\sum_i S_i+9\sum_i C_i-\sum_i V_i$ orders safe
completion, completion and fewer violations lexicographically: $9>8$ and
$81>9\cdot8+8$. The coefficients encode this chosen hierarchy; they were not tuned.
Partial panels consume budget but cannot nominate a policy.

The PPO comparison trains the same head-coordinate actor with a separate critic
and stochastic raw logits; evaluation exports a deterministic actor. Its terminal
reward is $(81S_i+9C_i-V_i)/90$, without a dense progress reward. Training and
selection together receive exactly 40,000 actions per optimizer per repeat.
PPO also learns four exploration-scale parameters and the critic; the comparison
matches actor coordinates, information and interactions, not total training parameters.
The finite-panel ordering does not imply lexicographic population probabilities:
PPO maximizes expected fixed utility. Safety specification and preference design
remain, and reduced tuning is untested. Only one PPO configuration was tested.

\subsection{TwoRoom evidence}
The initial replication found CMA safe completion of 8/36, 6/36 and 11/36 across
three repeats, compared with native 3/36 and random search 4/36, 4/36 and 1/36.
A later representation/mechanism study selected head coordinates and evaluated on
24 new layouts with three planner seeds. It found 45/216 safe completions for CMA
and 47/216 for random search, against shared native 7/72. One untrained reference
outperformed its paired evolved adapter. These outcomes limit any conclusion drawn
from the initial positive pilot.

\begin{table}[t]
\caption{Matched-objective TwoRoom study. Native is shared within each host; the
other rows pool three learner repeats on the same 24 layouts and three planner
seeds. Episodes are not independent task samples. All episodes engaged.}
\label{tab:ppo}
\centering\small
\begin{tabular}{lrrrrr}
\toprule
Controller & Safe & Complete & Violated & Unfinished & $n$\\
\midrule
Native & 14 & 77 & 75 & 67 & 144\\
Untrained & 34 & 127 & 111 & 89 & 216\\
CMA-ES & 44 & 116 & 109 & 100 & 216\\
Random & 37 & 99 & 84 & 117 & 216\\
PPO & 20 & 110 & 116 & 106 & 216\\
\bottomrule
\end{tabular}
\end{table}

The separate matched PPO study used two complete optimizer comparisons on one
RTX 4090 and one on an L4. CMA ranked highest in aggregate (Table~\ref{tab:ppo})
and exceeded PPO in every repeat. A pre-specified crossed bootstrap resampling
learner repeats and layouts, with planner seeds retained together, gave CMA minus
PPO $+11.11$ percentage points, exploratory 95\% interval $[-6.94,+28.24]$;
CMA minus random was $+3.24$ points, interval $[-11.57,+18.52]$.
Three repeats and shared final layouts limit inference. This supports neither
broad evolutionary superiority nor a claim that PPO cannot learn the task.

The mechanism study also examined task-only ranking, removal of online hazard
geometry and recurrent-state resets. Their safe-completion totals were 29/216,
36/216 and 37/216, respectively, with mixed effects across repeats. Hazard relocation
changed 72/72 main-adapter action sequences and 0/72 hazard-blind sequences. This
establishes response to hazard information, not necessarily a safer response.

\subsection{PushT baseline and qualification}
A first distinct-scene screen excluded entire scene groups from an 8,000-episode
retrieval memory. Native completed 4/16 episodes, safely in 2/16, with 4/16 violations
and 12/16 unfinished. The prerequisite gate failed, so no parameter probes or
learning followed. That result is retained.

We then reproduced the released AP-CEM protocol on its 128 queries with its full
16,688-episode memory, published package versions and fixed seeds. The released
evaluator completed 88/128 and the existing wrapper 89/128; both engaged in all
128 episodes, with no initially successful queries. They shared 88 successes and
39 failures. Query episodes were excluded, but sibling demonstrations from the
same scene could remain. The study used RTX 4090 rather than H20 hardware and a
different Python patch release. Near-matching aggregate performance does not imply
bitwise equivalence.

The earlier screen and reproduction differ in memory, queries and runtime together.
Their difference does not isolate a causal explanation. The reproduced baseline
supports a new episode-excluded, familiar-scene qualification; it does not validate
novel-scene generalization. That qualification fixes eight development scenes and
twelve reserved candidate final scenes, with safe and risky witnesses and sixteen
pre-drawn head-coordinate probes. The completed 329-episode qualification passed its independent audit. Native
completed 12/16 episodes but violated in all 16; five scenes passed admission.
Across sixteen head-coordinate probes, safe completion on the admitted subset
ranged only from 0/10 to 1/10 and occurred on just one scene. Across all scenes,
probe completion ranged from 0/16 to 4/16, below native 12/16, despite engagement
in every episode. The useful-authority gate failed, so no matched optimizer
training followed and the twelve reserved cases remain untouched. This separates
baseline competence from useful control authority in the tested adapter space.

\section{Intended main study: Block Plucking}
Block Plucking uses frozen UNISafe components and task-native physical violations.
Here ``native'' denotes the bare anchored planner. The safety filter computes
shadow signals but never vetoes the planner's commands. This is a study built on
UNISafe components, not a reproduction of its released filtered-policy evaluation.
The primary endpoint also differs from LeWM: after initial benchmark success,
the physical success predicate must hold for all 40 additional controls (two
seconds), with no violation. Benchmark success, sustained safe success and any
violation during continuation are recorded separately. ``Unfinished'' here means
no initial benchmark success, including episodes that ended on a violation.
The current pilot selects among retrieved trajectory candidates around a frozen
world-model planner, rather than applying the LeWM recurrent action residual.
Its five linear weights modify candidate ranking using standardized retrieval
distance, imagined cost, mean arm magnitude, closing fraction and demonstration
progress. Neutral weights recover the bare planner. This implementation difference
is material to the scope of transfer.

The completed three-repeat pilot evaluated 48 fresh starts per controller.
Sustained safe success was 10/48 native, 12/48 untrained, 7/48 random search and
10/48 CMA. Violations were 28/48, 25/48, 30/48 and 26/48, respectively. CMA's
safe-completion differences from native were $+2,+2,-4$ across repeats. The pooled
pre-registered verdict was no signal. The pilot matched training episode counts,
whereas LeWM matches actual simulator actions; these are not identical budgets.
CMA and random search also used different training-panel allocation rules in this
pilot, which limits optimizer interpretation.

The next Block Plucking study independently changes training-panel breadth,
nominee selection, elite re-evaluation and frozen-system candidate scores. Its
selector expands from five to eight weights, adding imagined terminal minimum
critic value, predicted failure probability and ensemble disagreement along the
imagined action block. These are signals from the pretrained frozen system.
Its
first completed repeat (R1) passed an independent saved-record audit. On 16 final
starts, sustained safe completion was 2 native, 6 untrained, 6 random and 5 CMA;
violations were 9, 9, 6 and 10, respectively. Benchmark successes were 10, 10, 9
and 7, leaving 6, 6, 7 and 9 unfinished. CMA therefore received the frozen mixed
verdict: three more sustained completions than native, but one more violation.
Its final mean beat its last elite on the selection panel. Two repeats remain
active, so no pooled conclusion follows. Since the changes are bundled, a gain
would not identify its cause. Comparisons stay within host because equal seeds
produce different trajectories across hosts.

The R1 audit checks all 1,156 completed episode records, including four host-local
smoke episodes (two more than the original budget table), search rankings, nominees
and final panels. It reconstructs outcomes from saved predicates without replaying
the simulator. This run used 93,457 controls. Setup and discarded launches are
accounted separately; incomplete attempts provide only recorded lower bounds.
Three withdrawn-repeat smoke episodes are documented, but their raw resource
records have not been located. Complete study-wide resource use is not verified.

\section{Safety dial and reliable gains}
A risk limit is useful only when candidate policies offer different task/risk
choices. A zero-rollout diagnostic used complete training panels from the matched
PPO study, choosing maximum completion subject to empirical violation limits
$0,0.25,0.5,0.75,1$. It found distinct selected CMA and random policies in all three
repeats, but only one selected PPO policy per repeat. Missing targets remained
unavailable. These are eight reused selection cases, so their risk estimates
cannot certify deployment probabilities.

The prospective test locked 17 distinct selected actors, shared native and
three untrained references before 756 episodes on twelve previously observed
development layouts and three planner seeds. All 772 records, including 16
engineering episodes, passed audit. Of 45 target assignments, 17 were unavailable;
five of 28 available assignments exceeded the requested empirical risk limit.
Four of six curves selecting multiple policies had nondecreasing violations;
R0 CMA and R2 random search were nonmonotonic. PPO selected one actor per repeat
at limits 0.75 and 1, with no lower targets available. Such constant curves do not
establish a useful dial.

For example, R0 CMA violated in 18/36 episodes at limit 0.25 and 16/36 at 0.5.
R1 CMA offered a risk change from 4/36 to 11/36, but completion was only 11/36
and 12/36, leaving 25/36 and 24/36 unfinished. Every policy engaged in all 36
episodes. Repeated targets reuse the same actor evaluation, so assignments are
not independent trials. This is limited development evidence of policy choice,
not a calibrated-risk result. The library was trained with the original fixed
utility, not these risk limits. Its 360,000 acquisition actions remain charged
separately from the new 105,532 actions. No reduced-tuning claim follows.

A separate L4 development study tests population four against the existing eight
and matched random search in two repeats, each at 40,000 actions per arm. The
hypothesis is that more CMA updates can help at a small interaction budget.
Its advancement criterion requires gains over both controls in both repeats
without higher violations than population eight. No result is yet available.
Passing that criterion would motivate fresh confirmation, not establish a
reliable evolutionary advantage by itself.

\section{Limitations and remaining evidence}
The present monitors and inputs still require designed safety definitions.
TwoRoom hazards are constructed; PushT protects only the pusher; Block Plucking
uses different physical predicates and sustained completion. These settings support
mechanistic comparisons, not a universal human-values alignment claim. Familiar
scene retrieval and previously observed development layouts restrict transfer
claims. Small numbers of learner repeats, model-based numerical variability and
one ordinary-learning configuration restrict optimizer conclusions.

Reduced tuning remains a distinct experiment: compare a frozen adaptive method
against a specified manually tuned reference, count every development interaction
for each, and evaluate on untouched tasks under a predeclared practical comparison
criterion. Current fixed utility coefficients and post-diagnostic design choices
must remain disclosed. No such benefit is claimed here.

Active development and Block Plucking outcomes remain necessary for final claims.
Failed prerequisites stay in the evidence. Current results support useful TwoRoom
adaptation and baseline qualification; cross-environment gains, reduced tuning and
reliable evolutionary superiority remain unestablished.

\section*{Evidence map and draft status}
Rule43 accounts 9 and 12--17 index completed LeWM evidence;
\texttt{SafeBlock/PilotPooled.md} records the Block pilot. The active studies are
\texttt{Search2.md} and \texttt{EvolutionBudget\allowbreak Study.md}.
Final-result integration, broader literature coverage and author/supervisor review
remain before submission.

\begin{thebibliography}{14}
\bibitem{lewm} Lucas Maes, Quentin Le Lidec, Damien Scieur, Yann LeCun, and
Randall Balestriero. 2026. \href{https://arxiv.org/abs/2603.19312}{LeWorldModel:
Stable End-to-End Joint-Embedding Predictive Architecture from Pixels}.
\bibitem{ap} Xvyuan Liu, Jianjie Fang, Wei Wu, Chen Gao, and Yong Li. 2026.
\href{https://arxiv.org/abs/2609.30036}{Aim Short to Reach Far: Your Frozen World
Model Can Plan Better Than You Think}.
\bibitem{unisafe} Junwon Seo, Kensuke Nakamura, and Andrea Bajcsy. 2025.
\href{https://proceedings.mlr.press/v305/seo25a.html}{Uncertainty-aware Latent
Safety Filters for Avoiding Out-of-Distribution Failures}.
\emph{Proceedings of the 9th Conference on Robot Learning}, PMLR 305, 4442--4472.
\bibitem{cma} Nikolaus Hansen. 2016.
\href{https://arxiv.org/abs/1604.00772}{The CMA Evolution Strategy: A Tutorial}.
\bibitem{ppo} John Schulman, Filip Wolski, Prafulla Dhariwal, Alec Radford, and
Oleg Klimov. 2017. \href{https://arxiv.org/abs/1707.06347}{Proximal Policy
Optimization Algorithms}.
\bibitem{cpo} Joshua Achiam, David Held, Aviv Tamar, and Pieter Abbeel. 2017.
\href{https://proceedings.mlr.press/v70/achiam17a.html}{Constrained Policy Optimization}.
\emph{Proceedings of the 34th International
Conference on Machine Learning}, PMLR 70, 22--31.
\bibitem{budgeted} Nicolas Carrara, Edouard Leurent, Romain Laroche, Tanguy Urvoy,
Odalric-Ambrym Maillard, and Olivier Pietquin. 2019.
\href{https://papers.nips.cc/paper_files/paper/2019/hash/4fe5149039b52765bde64beb9f674940-Abstract.html}{Budgeted
Reinforcement Learning in Continuous State Space}.
\emph{Advances in Neural Information Processing Systems 32}.
\bibitem{deb} Kalyanmoy Deb. 2000.
\href{https://doi.org/10.1016/S0045-7825(99)00389-8}{An efficient constraint handling
method for genetic algorithms}. \emph{Computer Methods in Applied Mechanics and
Engineering} 186, 2--4, 311--338.
\bibitem{ranking} Thomas P. Runarsson and Xin Yao. 2000.
\href{https://doi.org/10.1109/4235.873238}{Stochastic ranking for constrained
evolutionary optimization}. \emph{IEEE Transactions on Evolutionary Computation}
4, 3, 284--294.
\bibitem{snes} Lenz Belzner and Martin Wirsing. 2021.
\href{https://doi.org/10.1016/j.scico.2021.102620}{Synthesizing safe policies under
probabilistic constraints with reinforcement learning and Bayesian model checking}.
\emph{Science of Computer Programming} 206, Article 102620.
\bibitem{sos} Enrico Marchesini, Davide Corsi, and Alessandro Farinelli. 2022.
\href{https://ojs.aaai.org/index.php/AAAI/article/view/20737}{Exploring Safer Behaviors
for Deep Reinforcement Learning}. \emph{Proceedings of the AAAI Conference on
Artificial Intelligence} 36, 7, 7701--7709.
\bibitem{moerl} Adam Callaghan, Karl Mason, and Patrick Mannion. 2025.
\href{https://doi.org/10.1016/j.neucom.2025.129991}{Extending Evolution-Guided
Policy Gradient Learning into the multi-objective domain}.
\emph{Neurocomputing} 636, Article 129991.
\bibitem{crmopo} Shangding Gu, Bilgehan Sel, Yuhao Ding, Lu Wang, Qingwei Lin,
Alois Knoll, and Ming Jin. 2024.
\href{https://arxiv.org/abs/2405.16390}{Safe and Balanced: A Framework for
Constrained Multi-Objective Reinforcement Learning}. arXiv:2405.16390.
\bibitem{gsrl} Weiqin Chen and Santiago Paternain. 2024.
\href{https://proceedings.mlr.press/v242/chen24b.html}{Generalized constraint for
probabilistic safe reinforcement learning}. \emph{Proceedings of the 6th Annual
Learning for Dynamics \& Control Conference}, PMLR 242, 1606--1618.
\end{thebibliography}

\end{document}
